\section{Open Questions and Next Work}
\label{sec:checkpoints-open-questions}

\subsection{Known Gaps}
\label{sec:checkpoints-open-questions-gaps}

Five open gaps are documented in the gaps/ directory and accepted as non-blocking to the program
ALIVE verdict. Each is recorded with a severity, a blocking status, and a remediation plan.

\textbf{GAP\_001: wasm4pm-compat Graduation Boundary Signal Implementation.} The
\texttt{Admission<T, W>} graduation trigger that signals when a compat-layer type should graduate
to the wasm4pm engine is not yet implemented in code. The research program has specified the
graduation surface in the graduation-surface-ledger.yaml and graduation-law.yaml documents, and the
GGEN\_ECOSYSTEM\_INTEL audit 4 (Graduation Boundary) passes against these specifications. However,
the executable signal that fires the graduation transition at runtime is deferred to Phase 12. This
gap does not block the current ALIVE verdict; it blocks the full GGEN\_ECOSYSTEM\_MANUFACTURING\_001
workflow.

\textbf{GAP\_002: OTel Trace Deserialization Rule Completeness.} The full OTel-to-OCEL translation
rule set is not yet finalized. The GGEN\_OTEL\_WEAVER\_PI\_ALIVE\_001 checkpoint establishes the
boundary axioms (telemetry is feedstock, conformance is court, schema diffs are not process drift),
and the BridgeRx component implements the Version 2 to Version 1 schema translation. However, the
complete rule set covering all attribute namespaces, type format variants, and edge cases in
collector pipeline configuration is an open research question. GGEN\_OTEL\_WEAVER\_PI\_ALIVE\_002
is pending this completion.

\textbf{GAP\_003: M\&A Deck Rendering Performance Optimization.} The ma-deck.tera template produces
board-admissible M\&A projection artifacts, but the rendering performance for large diligence
workbooks (30,639-character master framework document) has not been benchmarked or optimized.
This is a Phase 12 candidate.

\textbf{GAP\_004: Blue River Dam Orchestrator Fault Tolerance.} The Blue River Dam orchestrator
(629 lines, 5 of 5 tests passing) encodes LTL invariants for soundness, completeness, consistency,
conformance, and auditability. The Doctor component implements rollback and remediation protocols.
However, the fault tolerance behavior under network partition, partial state corruption, or
concurrent MAPE-K loop contention has not been tested or formally specified. This is a Phase 12
candidate.

\textbf{GAP\_005: Advanced LTL Verification Coverage Expansion.} The Blue River Dam orchestrator
encodes a specific set of LTL invariants. Broader LTL property coverage---covering liveness
properties, fairness constraints, and domain-specific SLA encodings---is an open Phase 13 research
question. The current implementation provides a foundation; the expansion requires additional formal
methods research.

\subsection{Deferred Gates}
\label{sec:checkpoints-open-questions-deferred-gates}

The GGEN\_ECOSYSTEM\_INTEL Audit 1 (No DTO Flattening) is currently failing. The violation is
specific: \texttt{to\_json\_string()} in \texttt{compat/src/manufacturing/mod.rs} and
\texttt{receipt\_json()} in \texttt{compat/src/manufacturing/traits.rs} are JSON serialization
methods that belong in the wasm4pm engine, not the compat layer. The remediation plan is stated in
GGEN\_ECOSYSTEM\_INTEL\_ALIVE\_001: move both methods to wasm4pm, remove DTO flattening from the
compat public API, re-run the boundary audits, and seal GGEN\_ECOSYSTEM\_INTEL\_ALIVE\_002. The
estimated timeline is 4 to 6 hours. Until this gate passes, the GGEN manufacturing pipeline is
halted for projection surface manufacture (TypeScript bindings, WASM ABI, Component Model WIT).

The OTel deserialization gate (GGEN\_OTEL\_WEAVER\_PI\_ALIVE\_002) is deferred pending finalization
of the complete OTel-to-OCEL translation rule set. The PARTIAL checkpoint
GGEN\_OTEL\_WEAVER\_PI\_PARTIAL\_001 records the design gate as passed; the ALIVE verdict requires
the full experimental suite to be confirmed. The estimated timeline is 2 to 3 hours, and no
blockers prevent immediate work.

\subsection{Research Risks}
\label{sec:checkpoints-open-questions-risks}

The principal research risk in the checkpoint architecture is the gap between machine-verified gate
evidence and assertion-based gate evidence. The current ALIVE gate assessment relies on file counts
asserted in Markdown documents rather than on CI-logged outputs of the verification bash script.
RESEARCH\_CRITERIA.md defines the verification script but does not record that it was run and its
output logged as an artifact. This means an auditor must re-run the script to independently confirm
the counts rather than examining a stored script output artifact. Future checkpoint generations
should capture and store the script output as a sealed receipt. [AUTHOR THESIS: this represents
a structural gap in the current evidence chain that should be addressed before Phase 13 research
relies on the checkpoint counts for formal theorem proof citations.]

A second research risk is that the checkpoint ledger covers research artifact density (file counts)
but does not currently cover research artifact quality. Gate 1 certifies that 33 doctrine files
exist; it does not certify that each doctrine file cites at least one primary source. The
RESEARCH\_CRITERIA.md does list source citation as an additional ALIVE requirement, but the
verification script does not check it programmatically. [INTERPRETATION: quality gates will require
a different mechanism---potentially a separate audit workflow that samples doctrine files and checks
citation presence---rather than a simple count check.]

\subsection{Next Receipted Work}
\label{sec:checkpoints-open-questions-next-work}

Three downstream workflows are authorized and ready by the program ALIVE verdict:

\textbf{GGEN\_ECOSYSTEM\_MANUFACTURING\_001} requires DTO remediation (GAP\_002 resolution) before
manufacturing can begin. Upon completion, it will manufacture TypeScript projection bindings via
Specta, WASM ABI conversions, and Component Model WIT schemas, and seal the
GGEN\_ECOSYSTEM\_MANUFACTURING\_ALIVE\_001 checkpoint with a BLAKE3 receipt chain.

\textbf{OTEL\_TRACE\_DESERIALIZATION\_001} can proceed immediately with no blockers. It will finalize
the OTel-to-OCEL translation rule set, confirm EventLog admission compliance, manufacture integration
test fixtures, and seal GGEN\_OTEL\_WEAVER\_PI\_ALIVE\_002.

\textbf{M\&A\_DECK\_MANUFACTURING\_PIPELINE\_001} can proceed immediately. It will manufacture the
board-ready M\&A projection deck, the detailed diligence workbook, and the financial modeling
spreadsheet against the 42 M\&A claim files, seal each artifact with BLAKE3, and deliver the slide-
to-receipt map for board presentation. All 42 M\&A claim files are complete; the Synthesis Director
is authorized to manufacture; board presentation can proceed upon completion.

The next checkpoint to be sealed after these three workflows complete is expected to be
GGEN\_ECOSYSTEM\_MANUFACTURING\_ALIVE\_001, followed by GGEN\_OTEL\_WEAVER\_PI\_ALIVE\_002, followed
by the first production M\&A deck delivery receipt. Each sealing will extend the immutable ledger
and authorize the next tier of downstream manufacturing work.
